
\section{Error Analysis of the Tuned Models}
\label{sec:erroran}
The tuned models fail differently, and the difference is instructive.
Table~\ref{tab:erroran} gives both confusion matrices on the 624-film
official test set (from \texttt{tuned\_results.json}).
\begin{table}[h]
\centering\small
\begin{tabular}{lccccccc}
\hline
Model & TN & FP & FN & TP & Acc & Sens & Spec \\
\hline
CNN core & 135 & 99 & 15 & 375 & 81.73 & 0.962 & 0.577 \\
RegionGCN & 140 & 94 & 5 & 385 & 84.13 & 0.987 & 0.598 \\
\hline
\end{tabular}
\caption{Tuned-model confusion matrices, official test set (234 normal /
390 pneumonia).}
\label{tab:erroran}
\end{table}
The CNN core is a high-sensitivity screener: 96.2\% sensitivity at the
cost of 99 false positives (specificity 57.7\%). The RegionGCN improves \emph{both} error
types at once --- 94 false positives vs 99, and only 5 missed pneumonias
vs 15 --- winning on accuracy, sensitivity and specificity
simultaneously. Neither model dominates; the graph model's edge
is calibration (AUC 0.9471 vs 0.9174), which is what a clinician-facing tool
actually consumes.
\subsection{Where the remaining errors live}
Cross-referencing the 99 CNN false positives against the census flags:
the test set was expert-adjudicated, so census flags do not apply there,
but the \emph{acquisition} forensics do --- the false-positive films
share the low-contrast/low-sharpness signature measured on the training
flags. The residual error mass is concentrated on the same equivocal
acquisition stratum the census identified in training labels: the model
and the label process struggle in the same place. This convergence
between the error surface and the label-noise surface is the paper's
quiet internal consistency check: two independent measurements (a
label-process audit and a model error analysis) localize the same
stratum.
